% PRJ93 literature review — restructured draft, 2026-08-01.
% Working copy under brain/. NOT pushed to Overleaf (gate 5).
%
% Changes against the Overleaf version are itemised in
% brain/ledger/litreview_corpus_judgement.md and the diff summary in
% brain/ledger/litreview_critique.md.
%
% No project result figure appears in this chapter by design: the review
% anticipates the dissertation's findings by argument, never by quoting a
% number, so SKILL.md T1 has no surface here to fail on.

This chapter is organised as a methodological pipeline that mirrors the
proposed system: an agent must first learn a venue's normal trading rhythm,
then satisfy itself that the rhythm has been scored and selected honestly, then
notice when the present departs from it, then decide whether the departure is
worth raising, and finally be judged on whether that decision was a good one.
Each section consumes the output of the one before it. Two of them are
deliberately parallel: Section~\ref{sec:rw-ruler} asks how a forecaster should
be judged, and Section~\ref{sec:rw-evaluation} asks the same question of the
agent built on top of it.

Four arguments are threaded through the whole, and this review takes a position
on each rather than surveying opinion. Operational rhythm can be borrowed across
venues, but the benefit is a property of the size of the group and not of the
act of sharing, so a three-venue estate should expect little of it. A deviation
is best defined as the breach of a calibrated band, whose guarantee degrades
gracefully rather than holding through a regime change. A learned rhythm
functions as the agent's memory, in a bounded sense argued for below rather
than assumed. And the worth of a proactive alert turns on the asymmetric human
cost of a false alarm, which is an empirical claim with empirical support and
is treated here as one.

Two features of the evidence base are stated once, here, rather than hedged
throughout. First, the frontier of proactive agency is young: of the sources
cited in this chapter, eleven are 2025--26 preprints that have not completed
peer review, and they are concentrated in
Sections~\ref{sec:rw-surfacing} and~\ref{sec:rw-evaluation}. Where an argument
in this dissertation rests on one of them, it is flagged at the point of use,
and no claim of contribution rests on unrefereed work alone. Second, this review
was written after the experimental work it introduces, and it would have been
easy to write a chapter that predicted only the results that succeeded. It does
not: where the literature anticipates a negative outcome in this estate, that
prediction is made explicitly, before the result is reported.

\section{From Decision Support to the Autonomous Managerial Agent}
\label{sec:rw-framing}

The system this dissertation builds is a manager-facing agent, and a
related-work chapter must first settle what kind of artefact that makes it.
\citet{gorry_framework_1971} drew the distinction that
still organises the field: a decision is \emph{structured} when Simon's
intelligence, design and choice phases can be formalised into rules, and
\emph{unstructured} when it resists formalisation and demands human judgement,
with decision-support systems positioned to assist managers on the
unstructured and semi-structured cases rather than to replace them. For five
decades the manager remained the decision-maker and the software remained the
support. \citet{kumar_agentic_2026} argue that agentic AI breaks
this arrangement. What separates an agent from earlier enterprise software,
on their account, is delegated autonomy: the capacity to initiate actions,
make decisions and coordinate activity toward a goal with little human
prompting, so that agency is distributed across people and machines rather
than held by the user alone. Read against Gorry, this is not a faster support
tool but a different category of actor, one that participates in the decision
instead of merely informing it.

The deployment record indicates the category is already populated, and that it
is populated faster than it is understood. An index of thirty widely used
agentic products \citep{staufer_2025_2026} documents systems acting with limited
oversight, yet also finds that developers disclose almost nothing about their
safety and evaluation, leaving no public information for 135 of 240 safety
fields. The gap is not peculiar to agents. Surveying published reports of
machine-learning deployments across industries and use cases,
\citet{paleyes_challenges_2022} map the difficulties practitioners report onto
the stages of a deployment workflow and find that practitioners face issues at
every stage of that process --- a catalogue assembled from case studies of
systems in use, not from benchmark results. Two consequences follow for this
work. Principled evaluation of a deployed agent is a contribution in its own
right rather than an afterthought, and the setting in which such an agent is
evaluated is itself a variable, because the problems that a deployment reports
are not the problems a benchmark reports.

Adopting the actor view has one further consequence that the rest of this
chapter develops. An agent that acts unasked must hold an internal model of what
is \emph{normal} for the venue it manages, because an unprompted action is only
warranted by a departure from expectation. Learning that model from short,
noisy trading records is the constraint, named here and resolved later, that
shadows the whole design.

\section{Learning the Rhythm: Forecasting Normality Under Data Scarcity}
\label{sec:rw-rhythm}

By a venue's \emph{rhythm} we mean its recurring pattern of trade, the
day-of-week and seasonal regularities against which any given day can be
called typical or surprising. The defining difficulty is scarcity: a single
hospitality venue offers months, not years, of clean history, and a venue
that has just opened offers none at all. This is the cold-start constraint,
and it rules out the comfortable assumption of much forecasting work that each
series carries enough of its own history to be fitted in isolation.

Within that constraint, the demand-forecasting literature delivers a result
that is easy to overlook: simple methods are hard to beat.
\citet{chae_value_2024}, studying a large restaurant chain, show that deep
models using only internal operational and calendar features match or
outperform machine-learning models given access to richer external
macroeconomic and health data, so the marginal value of additional data is
smaller than practitioners assume. \citet{hossain_comparative_2025}
reinforce the point from the opposite direction: although a multilayer
perceptron and random forest rank highest on their restaurant data, classical
exponential smoothing and \citet{croston_forecasting_1972}'s method outperform
a gradient-boosted baseline, which fails to capture seasonality without
hand-built features. \citet{schmidt_machine_2022} sharpen the same point on a
single restaurant's sales: across twenty-three models and two baselines, the
best one-day-ahead result is a kernel-ridge model at 19.6\% sMAPE. At the
one-week horizon the ordering reverses, every non-recurrent model degrading
past 20\% error while the best recurrent model holds 19.5\%, so the
competitiveness of simple models is horizon-dependent rather than general.
Taken together these studies establish a baseline expectation this work
inherits: on small, seasonal hospitality series, classical statistical and
linear models set a competitive bar that a more elaborate method must clear on
evidence rather than on architecture.

\subsection*{When does borrowing across series pay?}

The transfer route promises to relieve scarcity by importing structure learned
elsewhere, and the question this section must settle is not whether that is
possible but when it is worth doing. \citet{montero-manso_principles_2021}
give the sharpest available answer, and it is not the one the enthusiasm for
shared models suggests. Globality is not a similarity assumption: global
methods are, in their words, not more restrictive than local methods, since
both can produce the same forecasts without any assumptions about similarity of
the series. What differs is where the complexity budget is spent. The
complexity of local methods grows with the size of the set while it remains
constant for global methods, so that --- and the qualifier is the authors' own
--- \emph{in large datasets} a global algorithm can afford to be quite complex
and still benefit from better generalisation. The benefit of pooling is
therefore a statement about the number of series available to amortise a
shared function over, not about the act of sharing. An estate of three venues
is not a large dataset in this sense, and the honest prediction the principle
licenses is that cross-series learning will do little for such an estate. This
dissertation reports exactly that outcome, and reports it as the predicted one.

With that expectation set, the mechanisms are worth stating precisely.
Time-series foundation models pretrain on large heterogeneous corpora and then
forecast unseen series zero-shot, and they divide by representation. TimesFM
\citep{das_decoder-only_2024} is a patched decoder that predicts the next patch
autoregressively and reaches near-supervised zero-shot accuracy with only 200M
parameters. Chronos \citep{ansari_chronos_2024} takes the opposite stance,
scaling and quantising real values into a fixed token vocabulary so an
off-the-shelf language-model architecture can be trained on them, and reports
comparable or superior zero-shot accuracy across 42 datasets. A wider family
--- masked-encoder, quantile, masked-reconstruction and probabilistic or
industrial entrants
\citep{woo_unified_2024,liu_moirai_2026,goswami_moment_2024,rasul_lag-llama_2024,garza_timegpt-1_2024}
--- fills out the design space without changing the central premise that
pretraining on time series, not on language, is what transfers.

That premise needs a sharp qualification, and \citet{tan_are_2024} supply it.
In ablations on three popular methods that bolt a pretrained language model
onto a forecaster, they find that removing the language model or replacing it
with a basic attention layer leaves accuracy unchanged or better while cutting
training and inference time by up to three orders of magnitude, and that the
language model helps neither sequence modelling nor the few-shot regime. This
does not contradict TimesFM and Chronos so much as locate the value precisely:
gains come from time-series pretraining and patching, not from language-model
machinery. TimesFM reaches the same conclusion from the other direction,
improving on language-model prompting by more than 25\%
\citep{das_decoder-only_2024}. For a small hospitality estate the lesson is a
disciplined baseline ladder, classical models first, then a transfer-based
global model, with scepticism toward any component justified by language-model
provenance alone.

One entrant deserves separate treatment, because it is the model this work
eventually serves. Chronos-2 \citep{ansari_chronos-2_2025} departs from its
predecessor by alternating standard time attention with a \emph{group}
attention layer that aggregates across all series sharing a group identifier at
each patch index. A group may be a set of series sharing a source or metadata,
so that the model performs cross-series learning at inference time instead of
forecasting each series from its own history alone, and the authors single out
short histories and cold starts as the conditions under which that sharing
helps most. Read against \citet{montero-manso_principles_2021}, that claim is
conditional in a way the architecture paper does not emphasise: group attention
supplies the mechanism for pooling but cannot manufacture the group size that
makes pooling pay.

The same mechanism admits past-only and known-future covariates as group
members, which places Chronos-2 among the small number of pretrained
forecasters that can condition on exogenous information at all. TabPFN-TS
\citep{hoo_tables_2026}, an 11M-parameter tabular foundation model applied to
forecasting through temporal featurisation, is the other, and it inherits from
TabPFN \citep{hollmann_accurate_2025} a design explicitly aimed at datasets of
up to ten thousand rows. \citet{hertel_explainable_2026} evaluate both on load
forecasting and find them competitive with transformers trained from scratch on
domain-specific data while requiring no training data at all, and report that
Moirai, by contrast, does not effectively leverage future covariates in that
setting. A single hospitality estate, with three venues, one year of trade and
a handful of calendar and weather regressors, sits inside the regime these two
models were designed for rather than at the scale the transformer literature
usually assumes, and the covariate capability is the specific reason one of
them is served here.

What makes transfer usable under genuine cold-start is the ability to adapt at
inference without per-venue retraining. \citet{das_-context_2025} demonstrate
in-context fine-tuning, prompting a foundation model with several related
series so it adapts to a target distribution at inference time, and report up
to 25\% improvement while rivalling a model explicitly fine-tuned on the
target, with no gradient updates. \citet{zhou_context-driven_2025} attack the
zero-history case directly, retrieving the traffic of semantically similar
pages to forecast a new page that has no series of its own, which is the
structural analogue of forecasting a venue on its opening week from its
siblings. GPHT \citep{liu_generative_2024} shows the underlying mechanism,
pretraining on a mixed channel-independent corpus to cut error by roughly
5.75\% over training from scratch. Each of these is a demonstration on a corpus
of many series, which is consistent with, rather than an exception to, the
locality--globality principle above.

Estate structure can also be exploited through reconciliation. MinT
\citep{wickramasuriya_optimal_2019} adjusts independent forecasts so that
venue-level predictions sum coherently to the estate total, minimising the
trace of the reconciled error covariance among all linear reconciliations that
preserve unbiasedness, an optimality that holds only if the base forecasts are
themselves unbiased --- a condition Section~\ref{sec:rw-ruler} shows is not
innocuous on a series of mostly zeros. HiGP \citep{cini_graph-based_2024} folds
the same coherence into an end-to-end model by projecting predictions onto the
space of coherent forecasts. The review by
\citet{athanasopoulos_forecast_2024} maps the resulting design space. These
results jointly justify treating a sparse or new venue not as an isolated
forecasting failure but as a borrower of rhythm from its estate and its
analogues --- while setting the expectation, from
\citet{montero-manso_principles_2021}, that on an estate this small the loan
will be a small one.

\subsection*{Intermittent trade is a different object}

The literature above assumes demand arrives continuously. For a venue that
trades one or two days in a typical week it does not, and the distinction is
not cosmetic: it changes the model, the metric and the meaning of a good
forecast. \citet{croston_forecasting_1972} separated the problem into the size
of a demand when it occurs and the interval between occurrences, estimating
each separately rather than smoothing a series that is mostly zeros;
\citet{syntetos_accuracy_2005} showed Croston's estimator is biased and
supplied the correction that bears their name. The classification that decides
when these methods apply is due to \citet{syntetos_categorization_2005}, who
partition series by average inter-demand interval and the squared coefficient
of variation of demand sizes into smooth, erratic, intermittent and lumpy
regions. \citet{kostenko_note_2006} then showed the published cutoffs contain
arithmetic errors, correcting them to $p = 4/3$ and $v = 0.5$ in place of 1.32
and 0.49, and derived the selection rule $v > 2 - \tfrac{3}{2}p$ dividing the
plane diagonally between Croston and the bias-corrected estimator. The
correction is small in magnitude and decisive in effect for any series whose
interval falls between the two constants, as this dissertation's anchor venue
does.

A second tradition models the same structure as two decisions rather than one.
\citet{cragg_statistical_1971} proposed the two-part specification in which a
binary model governs whether the outcome is positive and a separate model
governs the amount conditional on occurrence, and \citet{mullahy_specification_1986}
developed the unrestricted hurdle form, estimating the occurrence and magnitude
components by separate maximum likelihood. The attraction for a booking-led
venue is direct: whether the venue trades tomorrow is a different question,
answerable from different information, from how much it takes if it does. The
limitation is equally direct, and is a limitation of data rather than of
method --- a hurdle model is only as good as the occurrence signal available to
it, and where that signal lives outside the dataset the two-part structure
buys nothing.

\subsection*{Weather, and the temptation of exogenous data}

Weather is the exogenous covariate a hospitality forecaster reaches for first,
and the load-forecasting literature is the place where that reflex has been
tested most carefully. The result is a caution.
\citet{haben_short_2019} forecast demand at the low-voltage level and report
that, in contrast to high-voltage load forecasts, temperature was not an
important factor in the accuracy of their forecasts: often it had little to no
effect, and in many cases including temperature was actually detrimental to
accuracy. Their proposed explanation is the one that transfers: temperature is
strongly correlated with seasonality, so a model can learn temperature as a
surrogate for the seasonal pattern it already has, and produce larger errors
than it would with no temperature at all. The relevant variable is the level of
aggregation. Weather drives aggregate demand across a region; it competes with
the day-of-week pattern at the level of a single site.
\citet{hertel_explainable_2026} put a magnitude on the residual contribution,
attributing 3.55\% of a covariate-informed Chronos-2 forecast's feature
importance to temperature and 2.74\% to irradiance, against roughly 89\%
carried by the load history itself. A single hospitality venue is a site, not a
region, and the expectation this literature sets is that a weather covariate
will be marginal at best and capable of harm at worst. This dissertation's
weather result is consistent with that expectation, and the expectation is
recorded here rather than after the fact.

\subsection*{From a point to a band}

A point forecast, however accurate, cannot say what would count as surprising.
The conformal-prediction line replaces the point with an interval that carries
a distribution-free coverage guarantee. In its split form, which is the
variant used here, a score function is calibrated on held-out data and an
empirical quantile of those scores sets the interval width.
\citet{angelopoulos_conformal_2023} state the resulting guarantee in two
directions at once: coverage is bounded below by the nominal level and above
by that level plus a term of order one over the calibration-set size. The
upper bound deserves as much attention as the lower one, because it means a
band that contains every observation is not a cautious band but a
mis-calibrated one, and coverage reported without the width that produced it
conceals the difference.

That guarantee holds under exchangeability, which trading data has no
particular reason to satisfy, and the online variants exist to address exactly
this. Adaptive conformal inference \citep{gibbs_adaptive_2021} adjusts the
nominal level in response to realised coverage. EnbPI \citep{xu_conformal_2021}
discards the exchangeability requirement altogether, guaranteeing approximately
valid marginal coverage under mixing conditions on the error process rather
than distributional assumptions on the data, and calibrating by updating past
residuals through a sliding window without refitting the underlying model; its
successor \citep{xu_sequential_2023} makes the residual model explicit, and
conformal PID control \citep{angelopoulos_conformal_2023-1} recasts the
adjustment as feedback control. The practical guidance is collected by
\citet{stocker_gentle_2025}. Adaptivity carries a price that
\citet{zaffran_adaptive_2022} make explicit: on exchangeable scores the
adaptive update degrades efficiency in proportion to its step size, so a
method that adapts to a shift that never arrives is worse than one that does
not adapt at all. \emph{Calibration}, in this sense, is the property that a
stated 90\% interval contains the truth about 90\% of the time, neither more
nor less. Expressing a venue's rhythm as a calibrated band, rather than a
single expected number, is the design choice this section ends on, because once
the rhythm is an interval the question shifts from \emph{what will happen} to
\emph{what would count as surprising}.

\section{Scoring the Rhythm: Error Measures and Model Selection}
\label{sec:rw-ruler}

A rhythm that cannot be scored honestly cannot be defended, and two distinct
questions have to be answered before any forecast in this dissertation means
anything: which error measure expresses what a venue actually cares about, and
by what procedure one model is declared better than another. Both are usually
treated as preliminaries. On intermittent hospitality data they are the
substance.

The scaled measures are the natural candidates on series of different
magnitudes, but \citet{hewamalage_forecast_2023} show that the choice is not
neutral: measures built on absolute errors optimise for the median, and on
intermittent series that property makes a constant zero look like the best
available prediction. The same paper notes the companion hazard for
benchmark-scaled measures, that a naive benchmark scoring exact zeros on zero
actuals deflates the denominator those measures divide by. Both failure modes
bite hardest on precisely the venue this estate most needs a metric for --- one
whose median trading day takes nothing. The M5 competition
\citep{makridakis_m5_2022} answered this by scoring on the root mean squared
scaled error across a retail corpus in which 77.3\% of series had an average
inter-demand interval above $4/3$; squared base errors optimise for the mean
rather than the median \citep{hewamalage_forecast_2023}.
\citet{kolassa_why_2020} argues more generally that the measure encodes an
implicit choice of which functional of the predictive distribution is wanted
--- naming the mean, the median and the $(-1)$-median as the summaries
different measures reward --- so that selecting a measure is selecting a target
and not merely a scale.

Three further results converge on the same conclusion from different
directions, and together they make the choice of a squared scaled error less a
preference than a constraint. First, coherence:
\citet{kolassa_we_2023} states the cost in his title, though the
incompatibility is narrower, and more useful, than the title suggests. Because
the median of a sum is not generally the sum of the medians, point forecasts
minimising MAE, or MASE, which is a scaled MAE, are usually not coherent. The
exception is the instructive part --- the only error measures whose minimising
point forecasts are coherent are the squared error and monotonic functions of
weighted sums of squared errors, because the expectation is additive. The
conflict is a property of absolute-error measures, not of error minimisation as
such, and a squared-error measure dissolves it rather than trading against it.
For an estate that reconciles venue forecasts to an estate total, this is not
an aesthetic point: it is the condition under which the reconciliation of
Section~\ref{sec:rw-rhythm} is coherent with the objective the forecaster was
fitted to, and it is a second reason, independent of intermittency, to prefer
squared error.

Second, the limit of error as an instrument at all.
\citet{chatfield_all-zero_2007} show in a factorial study of lumpy demand that
the lowest forecasting error does not necessarily lead to the lowest system
cost, and --- the sharper half of the finding --- that all-zero forecasts yield
the lowest cost when lumpiness is high, and remain best at middling lumpiness
whenever shortage cost greatly exceeds holding cost. Where a series is lumpy
enough, error minimisation and cost minimisation come apart, and the honest
response is to change the objective rather than to keep tuning against a
measure that has stopped tracking what matters.

Third, the point forecast may be the wrong object.
\citet{kolassa_evaluating_2016} argues that the emphasis on point forecasts
misses the mark for count and intermittent demand, and that one should instead
characterise the entire predictive distribution, from which point, interval or
quantile forecasts can be extracted as required. His instrument is the proper
scoring rule, which assesses calibration and sharpness together rather than
either alone. This closes a loop opened at the end of
Section~\ref{sec:rw-rhythm}: a band reported by its coverage alone is scored on
calibration with sharpness suppressed, which is exactly the asymmetry
\citet{angelopoulos_conformal_2023}'s two-sided bound warns against. Coverage
and width are one measurement, not two, and a proper scoring rule is the
literature's name for treating them that way.

The second question is comparison. Declaring one model better than another on a
difference of means is not a finding, and the forecasting literature has said
so for thirty years. \citet{diebold_comparing_1995} formalised the test of
equal predictive accuracy between two forecast sequences;
\citet{harvey_testing_1997} showed the statistic is seriously over-sized in
moderate samples and supplied the finite-sample correction
$\left[\frac{n + 1 - 2h + n^{-1}h(h-1)}{n}\right]^{1/2}$, whose behaviour is
itself informative --- for a short evaluation window at a long horizon the
bracketed quantity can vanish, and where it does no corrected variant of the
test exists to be computed. That is a fact about the design rather than a
defect of the correction, and this dissertation states where it applies rather
than working around it. \citet{hansen_model_2011} generalise from the pairwise
case to the set: a model confidence set contains the best model with a given
level of confidence, and, in their own formulation, acknowledges the limitations
of the data, such that uninformative data yield a set with many models whereas
informative data yield a set with only a few. A wide retained set is therefore
a statement about the evidence rather than a failure of the procedure --- and
on six folds of a single year's trade, a wide set is what an honest procedure
should be expected to return.

Two recent studies show how fragile rankings are when this discipline is
absent. \citet{brigato_there_2025}, evaluating eight models on fourteen
datasets across roughly five thousand trained networks, find that slight
changes to experimental setups or evaluation metrics drastically shift the
common belief that newly published results advance the state of the art, and
that even minor deviations in setup can dramatically change performance
rankings. \citet{hewamalage_look_2021} make the same concern measurable,
proposing rank stability --- how much the ranking of an experiment changes
across similar datasets when models and errors are held constant --- and
finding that the main drivers of instability are hierarchical aggregation and
scaling, with the scale normalisation of the M5 error measure less stable than
other scale-free errors. Two things follow. The number of evaluation origins is
not a free parameter, because a ranking computed on too few can reverse when
more are used; and the case for a squared scaled error assembled above is
strong but not costless, since the scale normalisation it depends on is itself
a documented source of rank instability. Both cautions are taken up in the
methodology, and both anticipate a reversal this dissertation reports.

\section{Noticing Deviation: From Calibrated Forecast to Signal}
\label{sec:rw-deviation}

Once the rhythm is a calibrated interval rather than a point, a meaningful
\emph{deviation} is definitionally a breach of that band, and the detection
literature's job is to make breach-detection reliable and honestly evaluated
on short, noisy series. The classical methods organise by what they detect.
Bayesian online change-point detection \citep{adams_bayesian_2007} maintains a
posterior over the run length since the last change, suited to abrupt regime
shifts such as a venue's permanent closure; the offline survey by
\citet{truong_selective_2020} reduces the whole family to a choice of cost
function, search method and penalty, which clarifies that change-point
detection answers \emph{when the regime moved}. That reduction also locates
the cumulative-sum scheme of \citet{page_continuous_1954}, the sequential
counterpart to the offline cost functions the survey enumerates. Page assigns
a score to each observation, accumulates those scores, and signals when the
running total has risen a fixed amount above its own running minimum, which
makes the scheme sensitive to small sustained shifts that a per-observation
test will not register. He also supplies the criterion by which such a
detector should be judged. The average run length, which he defines as the
expected number of observations before action is taken, measures the expense
of false alarms when conditions are stable and the delay before a real change
is acted on when they are not. That two-sided cost is the same one this
chapter arrives at independently in Section~\ref{sec:rw-evaluation}, stated
for an industrial inspection scheme seventy years before it was restated for
conversational agents. Being cheap to evaluate on each new observation and
requiring no posterior over run length, the cumulative-sum scheme is the
detector this work puts into production, with Bayesian online change-point
detection retained as the offline benchmark it is scored against.
Point-anomaly methods answer a different question, \emph{is this observation
extreme}, and SPOT \citep{siffer_anomaly_2017} answers it by fitting a
generalised Pareto tail via extreme value theory, requiring only a risk level
and assuming nothing about the data distribution. Mapping these onto a venue is
straightforward: a sudden quiet Saturday is a point anomaly against the band, a
venue closure is a change point, and conflating the two is a modelling error.

The harder problem is evaluation, and here the literature contains a finding
that should unsettle any reported detection score. \citet{kim_towards_2022}
show that the point-adjustment protocol, which counts a whole anomalous segment
as detected if any single point inside it crosses the threshold, is so generous
that randomly generated anomaly scores reach an adjusted F1 near one and
overturn most state-of-the-art results. \citet{liu_elephant_2024} confirm and
extend this with TSB-AD, a curated benchmark of 1070 series, identifying VUS-PR
as the measure that stays stable under small detection lags while remaining
unbiased against random scores, and finding along the way that simpler
statistical detectors often beat elaborate neural ones, which echoes the
baseline lesson of Section~\ref{sec:rw-rhythm}. Corrective protocols exist that
penalise the false positives the standard adjustment ignores, whether by
balancing the adjustment or by decaying credit for late detection
\citep{bhattacharya_towards_2024,gim_evaluation_2023}. The implication for this
work is methodological rather than incidental: detection on hospitality data
should be reported with a lag-tolerant, random-robust measure in the spirit of
VUS-PR, and point-adjusted F1 is rejected as a headline metric because it
would flatter a detector that has learned nothing.

The thread of Section~\ref{sec:rw-rhythm} closes here. CPTC
\citep{sun_conformal_2025} fuses the two halves directly, pairing a switching
dynamical model that predicts the underlying state with online conformal
prediction calibrated separately per state. Its guarantees repay precise
statement, because they are narrower than the framing invites. Exact
finite-sample coverage holds only under exchangeability, which a change point
violates by construction. What survives a regime shift is time-averaged
asymptotic validity, resting on an assumption that the states themselves have
a stationary distribution, together with faster post-shift convergence than
adaptive conformal inference, which the authors describe as shortening the
miscoverage period rather than removing it. The residual risk then concentrates
in a single quantity, since coverage degrades in proportion to the rate at
which the state is misclassified. \citet{barber_conformal_2023} give the
general form of the same trade, bounding the loss of coverage under arbitrary
departures from exchangeability by a weighted sum of total-variation distances
between residual vectors, and doing so with no assumption whatever on the joint
distribution of the data.

Read together, these results reposition the design question this chapter has
been building toward. A band cannot be promised to hold through a regime
change. It can only be promised to lose coverage in proportion to how badly
the regime is estimated. For a hospitality estate the operational consequence
is to prefer a regime variable that is observed over one that is inferred,
because whether a venue is trading is recorded in the till rather than hidden
in the residuals, and an observed state drives the misclassification term to
zero. A detector can establish \emph{that} trade has broken from its rhythm;
it cannot decide whether the break is worth a manager's attention.

\section{Reasoning and Surfacing: Agents that Act Without Being Asked}
\label{sec:rw-surfacing}

Converting a deviation signal into a timely, well-judged intervention needs
more than tool-use plumbing; it needs an agent that initiates action and
decides what is worth raising. The plumbing itself is mature: interleaved
reasoning and tool calls, self-supervised API invocation, and verbal
self-correction across attempts are established primitives
\citep{yao_react_2022,schick_toolformer_2023,shinn_reflexion_2023}, they are
present in the existing product, and they are not where the open problem lies.
That the problem is open is visible in evaluation: $\tau$-bench
\citep{yao_-bench_2024} places agents in multi-turn tool-and-user interactions
and finds frontier models far from solving them, with even GPT-4o completing
only 35.2\% of tasks in the airline domain, so reliability rather than
capability is the binding constraint even before proactivity enters.

Memory is the bridge from the forecasting half of this work to the agent half,
and the connection must be stated carefully, because it is easy to claim more
for it than is warranted. Retrieval-augmented generation
\citep{lewis_retrieval-augmented_2020} established the template of grounding a
generator in an external, non-parametric store accessed by a learned
retriever. \citet{hu_memory_2026} argue that agent memory is categorically more
than this: where classical retrieval reads from a static corpus per query,
agent memory is a persistent, self-evolving cognitive state that accumulates
the agent's own experience over time. Generative agents
\citep{park_generative_2023} make the mechanism concrete, with a memory stream
surfaced by a weighted combination of recency, relevance and importance, and
recursive reflection --- triggered when accumulated importance crosses a
threshold, rather than on a clock --- that abstracts raw records into
higher-level inferences.

The reframing this work takes from that literature is deliberately bounded. A
learned rhythm is a durable, evolving record of normality that the agent
retrieves against when deciding whether the present is worth remarking on, and
in that specific respect it does the work of the retrieval half of a memory
stream: an importance-triggered mechanism firing on a threshold is a closer
analogue to a band breach than a clock is. What a forecasting band does not
supply is the rest of the architecture --- there is no reflection step
abstracting episodes into higher-level inferences, and no store of the agent's
own past interventions to retrieve from. The claim made here, and defended in
the discussion, is the narrow one: the rhythm serves as the agent's model of
normality, not as a full memory system in the sense \citet{hu_memory_2026} and
\citet{park_generative_2023} describe. Retrieval over live operational data
carries a caveat in either case, since such stores are an attack surface, as
PoisonedRAG \citep{zou_poisonedrag_2025} demonstrates, which matters once the
memory is populated from production traffic.

The frontier of proactive agency is where this work positions its contribution,
and it is candid about its empirical maturity: most of the results in this
paragraph are preprints, and their claims are reported with corresponding
caution. \citet{lu_proactive_2024} provide the template. Their ProactiveBench
gathers 6,790 training and 233 test events from real human activity, trains an
agent to anticipate needs, and, tellingly, the dominant failure they document
is over-offering: even strong proprietary models carry false-alarm rates above
50\%, predicting help when no clear intent exists, so their best fine-tuned
agent reaches only a 66.47\% F1 while a separately trained reward model attains
91.80\% agreement with human judges. The headline finding is therefore not that
proactivity is hard to produce but that it is hard to \emph{restrain}.
Subsequent work attacks restraint from different angles. ProActor
\citep{ding_proactor_2026} reframes timing as a window of valid moments rather
than a single labelled point and optimises it with turn-level reinforcement
learning, on the argument that penalising any deviation from one rigid
timestamp is the wrong objective. ProAgentBench \citep{tang_proagentbench_2026}
grounds the question in 28,000 events from over 500 hours of real sessions and
states the cost structure plainly: low precision in deciding when to assist
causes alert fatigue, where excess false alarms drive users to ignore or
disable the feature. PRISM \citep{fu_prism_2026} draws the consequence,
formulating intervention as a cost-sensitive decision in which the agent speaks
only when calibrated acceptance probability clears a threshold set by the
asymmetric costs of false alarms and misses, and reports a 22.78\% reduction in
false alarms with a 20.14\% F1 gain. \citet{gulati_ask_2026} add that the
urgency of a clarifying question depends on what is missing: goal clarification
decays quickly, retaining near-oracle value only if asked very early, while
input clarification remains recoverable through roughly half of execution, so a
single uniform urgency for every kind of missing information is wrong. Adjacent
systems for dialogue, wearables and mobile interfaces
\citep{liu_proactiveeval_2025,yang_contextagent_2025,yang_fingertip_2025}
confirm proactivity is being pursued across domains without supplying a
precedent in a multi-venue operational setting. Read together, this body of
work converges on one diagnosis that the next section takes up: the binding
failure of a proactive agent is the false alarm, and the open question is how
to weigh it.

\section{Judging the Judgement: Evaluating Proactive Intervention}
\label{sec:rw-evaluation}

The agent's decision to speak is only as good as our ability to tell a useful
prompt from an unwelcome one, and accuracy is the wrong instrument for that
measurement. Two literatures supply the right one. The first concerns
automated evaluation. \citet{zheng_judging_2023} show that a strong language
model as judge agrees with human preference above 80\%; on MT-bench pairwise
comparison with ties excluded, GPT-4 reaches 85\% agreement with humans against
81\% agreement among humans themselves, which makes the approach attractive at
scale. The qualifications are serious, however. \citet{wang_large_2024}
demonstrate a positional bias severe enough that swapping the order of two
candidates can reverse the verdict; \citet{panickssery_llm_2024} show that
judges recognise and favour their own generations, with self-preference scaling
with self-recognition; and JUDGE-BENCH \citep{bavaresco_llms_2025} finds
substantial variance in agreement across models and datasets, with agreement
depending on the property being evaluated, the expertise of the human judges,
and whether the text being judged is human- or model-generated, concluding that
judges must be carefully validated against human judgments before they are used
as evaluators. The position this work takes is to use an automated judge only
as an instrumented, bias-audited proxy, with the real target being human
adopt-or-dismiss outcomes on actual surfaced insights.

The second literature defines what the metric must encode, and it is the
human-factors account of trust in automation. \citet{meyer_conceptual_2004}
draws the distinction that sets the loss function: \emph{compliance} is the
operator's response to a warning, \emph{reliance} the response to its absence,
and crucially the two are governed by different properties of the detector.
Compliance depends on the response criterion, responses to warnings becoming
less likely when the probability of an unjustified warning is high, the
cry-wolf syndrome in which operators cease or slow their response. Reliance
depends instead on the system's sensitivity, and therefore on what it fails to
catch.

That the two failures are not merely different but unequal is an empirical
claim, and it has been tested directly. \citet{dixon_independence_2007} placed
operators under diagnostic automation that was either false-alarm-prone or
miss-prone and found that false-alarm-prone automation hurt overall performance
more than miss-prone automation, and that it degraded both compliance and
reliance while miss-prone automation appeared to affect only reliance. Their
stated implication for design is the premise this dissertation's loss function
encodes: false alarms are more damaging to overall performance than misses. The
evidence is a laboratory study of thirty-two participants on a synthetic
monitoring task, and the asymmetry is asserted here at that strength and no
higher; the authors themselves conclude only that compliance and reliance are
not entirely independent.

Field evidence puts a magnitude on the same mechanism.
\citet{ancker_effects_2017}, studying 112 clinicians over three and a half
years of electronic health record data, found that the likelihood of accepting
a reminder dropped by 30\% for each additional reminder received in an
encounter, and by 10\% for each five-percentage-point increase in the
proportion of reminders that were repeats. The setting is clinical rather than
hospitality and the number does not transfer, but the mechanism does, and it is
the mechanism a proactive agent must budget against: the cost of an
unnecessary alert is not paid by that alert, it is paid by the next one. The
structure is not new, and it is worth noting that the detection literature
reached it first --- Page's average run length \citep{page_continuous_1954}
already traded the expense of false alarms under stable conditions against the
delay before a real change is acted on. What the human-factors work adds is
that in a managerial setting the two costs are borne by a person whose
willingness to keep reading is itself consumed by the false alarms.

The remaining literature situates that asymmetry in a theory of appropriate
use. \citet{lee_trust_2004} frame the goal as appropriate reliance, calibrated
trust in which trust matches capability, with overtrust producing misuse and
distrust producing disuse, the inappropriate rejection of a capable aid that
\citet{parasuraman_humans_1997} tie directly to high false-alarm rates.
\citet{parasuraman_complacency_2010} document the opposite pole, automation
bias and complacency, and a meta-analysis \citep{hancock_meta-analysis_2011}
shows that system performance, not surface attributes, is the largest driver of
trust, with robot performance factors ($\bar{r} = +0.34$) far exceeding robot
attributes ($\bar{r} = +0.03$).

The consequence for evaluation is that the right metric is calibrated precision
against real dismissals, weighted by the asymmetric cost of an interruption,
not an undifferentiated F1. A near-ready template exists: HiL-Bench
\citep{trinh_hil-bench_2026} scores agents with Ask-F1, the harmonic mean of
question precision and blocker recall, which penalises both over-asking and
silent wrong guessing; adapted from its coding domain to adopt-versus-dismiss
outcomes, it captures exactly the two-sided failure a hospitality agent must
avoid. Calibration closes the loop back to the forecasting half:
\citet{guo_calibration_2017} show modern networks are overconfident and that
temperature scaling restores calibration measured by expected calibration
error, so the coverage guarantee that defined a deviation in
Section~\ref{sec:rw-deviation}, and the calibration-with-sharpness argument of
Section~\ref{sec:rw-ruler}, reappear here as an evaluation guarantee on the
agent's stated confidence.

\section{Synthesis: What the Literature Leaves Open}
\label{sec:rw-synthesis}

Four results stand established by the work above, and each is established with
a qualification this chapter has argued for rather than inherited. Operational
rhythm can be borrowed across venues and analogues despite scarce data
\citep{das_-context_2025,zhou_context-driven_2025,liu_generative_2024}, but the
benefit scales with the number of series available to pool over, so a small
estate should expect a small loan \citep{montero-manso_principles_2021}. A
deviation can be defined as the breach of a calibrated band whose loss of
coverage under a regime change is bounded rather than eliminated
\citep{sun_conformal_2025,barber_conformal_2023,liu_elephant_2024}. An agent
can be made to act unprompted, with its learned rhythm serving as the model of
normality it retrieves against, though not as a full memory architecture
\citep{lu_proactive_2024,hu_memory_2026,park_generative_2023}. And the worth of
such an action is measurable against the asymmetric cost of a false alarm, an
asymmetry that is empirically tested rather than merely assumed
\citep{meyer_conceptual_2004,dixon_independence_2007,ancker_effects_2017,fu_prism_2026,trinh_hil-bench_2026}.
Each has been demonstrated in isolation, and in domains --- coding, dialogue,
web traffic, wearables, clinical informatics --- distant from hospitality
operations.

The fourth demands the most care, because the method it calls for already
exists. PRISM \citep{fu_prism_2026} sets its intervention threshold from an
asymmetric miss-to-false-alarm cost ratio, and among the proactive systems
surveyed here it is the one that gates on a calibrated acceptance probability
rather than on accuracy alone. It reports the gate's effect, a 22.78\%
reduction in false alarms and a 20.14\% gain in F1, but not the calibration of
the probability the gate depends on. No claim of methodological novelty in
cost-sensitive intervention is therefore available to this dissertation, and
none is made. What PRISM leaves undone, and what none of the surveyed systems
attempts, is to fix that cost ratio to the stated preferences of the person who
bears the cost, to measure the calibration of the acceptance probability that
ratio is applied to, and then to score the resulting decisions against that
same person's accept-or-dismiss judgements on signals drawn from their own
business. Every system reviewed above is evaluated against an annotated corpus,
a simulated user, or a language model acting as judge.

The contribution is accordingly one of field instantiation rather than of
method, and that is a recognised category rather than a consolation.
\citet{paleyes_challenges_2022}'s survey is assembled from reports of systems
in deployment and finds practitioners encountering difficulties at every stage
of the deployment workflow --- a body of knowledge that exists because
deployments were written up, not because benchmarks were run. The contribution
here is to take a risk-sensitive intervention policy of the kind PRISM
specifies, place it above a rhythm borrowed across a small real estate and a
deviation signal defined by a calibrated band, and evaluate the assembly
against an operator's own decisions in a live multi-venue setting.

What this chapter has additionally tried to do is state in advance where that
assembly should be expected to fail. The locality--globality principle predicts
that pooling across three venues will buy little
\citep{montero-manso_principles_2021}. The low-voltage load-forecasting
evidence predicts that weather will be marginal and may be harmful
\citep{haben_short_2019,hertel_explainable_2026}. The rank-stability and
benchmarking literature predicts that a model ranking computed on a handful of
evaluation origins can reverse when more are used
\citep{hewamalage_look_2021,brigato_there_2025}. Each of those predictions is
borne out in the chapters that follow. Whether a policy developed on annotated
benchmarks survives contact with one manager, three venues and a year of untidy
till data is the remaining empirical question the literature reviewed here
leaves open, and it is the question this dissertation sets out to answer.
